% =========================================================================
% CONCLUSIONS — ready to paste
%
% The thesis file currently has, near line 2486:
%
%     \unnumberedsection{Conclusions}
%     \label{chap:conclusions}
%
% followed immediately by the bibliography. The regulation lists project
% conclusions as a mandatory element (item j), so this section cannot be left
% empty. Paste the text below between the \label and the bibliography.
%
% Written to stand on the results already in the thesis. If the three-seed
% comparison of the two centre-update rules completes, add the marked
% paragraph; if it does not, the section is complete without it.
% =========================================================================

\unnumberedsection{Conclusions}
\label{chap:conclusions}

This work set out to recognise six categories of indoor scene with a
transformer trained from scratch, to add a clustering objective to the
embedding space, to exploit unlabelled images through self-labelling, and to
analyse the outcome with classification metrics and representation-space
measurements. Each of the four was carried out, and the results are more
informative about the interaction between them than about any one in isolation.

The clearest result is that the choice of feature extractor dominated
everything else that was varied. A ResNet-18 of eleven million parameters
reached \(81.83\) percent on the test partition where a ViT-S/16 of
twenty-two million, trained on the same images under the same schedule for the
same number of epochs, reached \(56.33\) percent. No auxiliary objective, no
selection rule and no coefficient examined in this work moved accuracy by more
than a few percentage points. On a dataset of nineteen thousand images, the
inductive bias of the architecture is the decisive factor, and a transformer
trained from scratch at this scale does not have enough data to compensate for
lacking one.

Center Loss did what its formulation promises, and the promise turned out to be
narrower than expected. On MNIST, where the validation network already
separates the classes and reaches \(98.39\) percent, adding the objective
reduced the mean within-class distance from \(50.16\) to \(0.36\) and raised the
ratio of between-class to within-class distance from \(4.76\) to \(20.76\); the
two-dimensional embeddings show the tight, radially separated clusters reported
in the original work. That result establishes that the implementation is
correct, which is the necessary precondition for interpreting anything that
follows.

On the room dataset the same objective contracted the representation just as
decisively and organised it not at all. The mean within-class distance fell from
\(8.203\) to \(3.814\) for the transformer, a reduction of fifty-four percent --
and the mean distance between class centroids fell from \(7.907\) to \(3.389\),
a reduction of fifty-seven percent. Because both quantities shrank in the same
proportion, the scale-invariant measures did not improve: the separation ratio
fell from \(0.964\) to \(0.888\) and the silhouette coefficient from \(0.036\) to
\(0.026\). The convolutional model, whose separation ratio already exceeded one,
responded in the opposite direction, its ratio rising from \(1.465\) to
\(1.606\). Contraction and separation are therefore not the same thing, and an
objective that produces the first does not necessarily produce the second.

The pattern across the three settings suggests that the geometric consequence of
Center Loss depends on the structure already present in the representation it is
applied to. Where the classes are distinct, pulling samples towards their
centres tightens clusters that are already apart. Where the classes overlap in
the feature space, as six indoor scene categories with shared objects and shared
layouts largely do, the same force contracts signal and noise alike. This is
offered as an interpretation consistent with three observations, not as a
demonstration; separating it from the alternative explanations would require
varying the representation quality deliberately rather than observing it across
architectures that differ in other ways as well.

The self-labelling experiments did not harm performance and did not clearly
improve it. The best configuration exceeded the supervised baseline by
\(1.17\) percentage points, which is seven images out of six hundred and below
the resolution of a single comparison. What was consistent across every
configuration and both partitions was the influence of the checkpoint the
procedure started from: runs warm-started from the Cross-Entropy model
outperformed those warm-started from the Center Loss model in every case. Self-
training amplifies the representation it begins with rather than repairing it,
which is the same conclusion the geometric measurements reach by another route.

One methodological result deserves to be stated on its own. Two runs of an
identical configuration, differing only in the order in which images were
presented and in the hardware used, produced test accuracies \(3.50\) percentage
points apart. That is larger than any difference this work attributes to a
method. It is the reason the conclusions above rest on representation geometry,
which is measured over six hundred embeddings and admits interval estimates,
rather than on accuracy, which on a test set of this size cannot resolve the
effects in question.

% ---------------------------------------------------------------------
% ADD ONLY IF THE THREE-SEED COMPARISON COMPLETES.
% Replace the bracketed values with the measured ones.
% ---------------------------------------------------------------------
% A final observation concerns the update rule itself. The variant used
% throughout this work treats the class centres as parameters of the combined
% objective, so the coefficient \(\lambda_{\mathrm{C}}\) scales their gradient;
% the rule published by Wen et al. moves them with a dedicated rate,
% independent of that coefficient. The two are not interchangeable in practice.
% Under the variant used here the centres remained close to their random
% initialisation, their mean norm moving from \(1.96\) to \(2.00\) over fifty
% epochs while the embeddings they were pulling towards had a scale four times
% larger; their cosine alignment with the true class centroids reached
% \(0.78\). Under the published rule the centres tracked the centroids almost
% exactly, with an alignment of \(0.997\) and a norm of \(14.3\). The
% contraction reported above is therefore best understood as motion towards
% points that never became class prototypes.
% ---------------------------------------------------------------------

Three directions follow from these results. The first is to establish how much
of what has been measured survives repetition: every configuration reported here
was trained once, and the run-to-run variation quantified above is large enough
that the smaller effects should be treated as descriptions of particular runs
until they are reproduced across several seeds. The second is to separate the
weight of the representation term from the rate at which the class centres move,
which the implementation used here couples and which the coefficient sweep
therefore cannot disentangle. The third is to test the interpretation offered
above directly, by applying the same objective to representations of
deliberately varied quality on a single architecture, rather than inferring it
from architectures that differ in several respects at once.
